\documentclass[12pt]{article}
\usepackage[T1]{fontenc}
\usepackage[utf8]{inputenc}
\usepackage{lmodern}
\usepackage{textcomp}
\usepackage{enumitem}
\usepackage{geometry}
\geometry{margin=1in}
\begin{document}
\begin{center}
Answer Key - Philosophy - Undergraduate Level Assessment 22 \
\small (Answers are ordered exactly as in the question paper.)
\end{center}

\section*{Multiple Choice}
\begin{enumerate}[label=\arabic*.]
\item \textbf{Answer:} A
\\ \textit{Options:} A) theory of existentialism.; B) version of utilitarianism.; C) theory of subjectivism.; D) Kantian theory.; E) version of virtue ethics.
\\ \textit{Note:} Locke's emphasis on individual experience and the formation of identity through personal choices aligns with existentialist themes, highlighting the importance of existence preceding essence.
\item \textbf{Answer:} A
\\ \textit{Options:} A) TlcL; B) cTL; C) lTc; D) Tll; E) cTl
\\ \textit{Note:} The correct answer is TlcL because it accurately represents the relationship stated in the question, where "T" indicates the taller-than relationship, with "l" representing Leo and "c" representing Cathy, thus forming the expression "Leo is taller than Cathy."
\item \textbf{Answer:} A
\\ \textit{Options:} A) Kannon; B) Shakyamuni; C) Manjusri; D) Guan-yin; E) Tara
\\ \textit{Note:} Kannon, also known as Avalokiteshvara, is recognized in Mahayana Buddhism as the bodhisattva who embodies compassion and is specifically venerated for guiding the souls of deceased children to a peaceful afterlife, thus fulfilling the role of a savior in the context of Buddhist belief.
\item \textbf{Answer:} A
\\ \textit{Options:} A) deontological ethics; B) rule-based hedonism; C) act-based deontology; D) egoistic hedonism; E) moral relativism
\\ \textit{Note:} Dershowitz argues that deontological ethics, which focuses on the morality of actions based on rules and duties rather than consequences alone, can justify acts of terrorism in extreme cases by prioritizing the obligation to protect innocent lives over the inherent wrongness of the act itself.
\item \textbf{Answer:} C
\\ \textit{Options:} A) encourage students to think pathologically.; B) ill-prepare them for the workforce.; C) harm their ability to learn.; D) all of the above.
\\ \textit{Note:} Lukianoff and Haidt argue that institutionalizing vindictive protectiveness undermines students' resilience and critical thinking skills, ultimately impairing their capacity to engage with diverse ideas and learn effectively.
\end{enumerate}

\section*{True / False}
\begin{enumerate}[label=\arabic*.]
\setcounter{enumi}{5}
\item \textbf{Answer:} False
\\ \textit{Options:} A) True; B) False
\\ \textit{Note:} The retributive theory posits that punishment should be proportionate to the wrongdoing and takes into account the nature of the crime and the intentions of the offender, rather than asserting that all wrongdoers deserve punishment indiscriminately.
\item \textbf{Answer:} False
\\ \textit{Options:} A) True; B) False
\\ \textit{Note:} Cicero argued that true morality aligns with Nature and that expediency, which often prioritizes self-interest over moral principles, is not in harmony with natural law.
\item \textbf{Answer:} False
\\ \textit{Options:} A) True; B) False
\\ \textit{Note:} The statement is false because successful interrogations can be achieved through ethical techniques such as building rapport and using psychological strategies, which are more effective and humane than torture.
\item \textbf{Answer:} True
\\ \textit{Options:} A) True; B) False
\\ \textit{Note:} Metz asserts that the essence of human dignity is rooted in our ability to engage in deep, meaningful relationships with others, emphasizing that our social connections and communal bonds are integral to our sense of worth and value.
\item \textbf{Answer:} False
\\ \textit{Options:} A) True; B) False
\\ \textit{Note:} The Babylonian captivity, also known as the Exile, actually began in 586 BCE when the Babylonians conquered Jerusalem, long before the rise of the Macedonian Empire.
\end{enumerate}

\section*{Long Form Response}
\begin{enumerate}[label=\arabic*.]
\setcounter{enumi}{10}
\item \textbf{Answer:} In John Rawls's concept of the original position, individuals are placed behind a "veil of ignorance," which strips them of any knowledge regarding their personal circumstances, including their social status, wealth, abilities, and personal values. This lack of information is crucial as it encourages decision-making based on principles of justice and fairness, rather than self-interest. By being unaware of their specific situations, individuals are motivated to establish rules that ensure equitable treatment for all, as they could potentially find themselves in any position within the society they are designing. Consequently, this framework fosters impartiality and prioritizes the creation of a just society, wherein the rights and opportunities of each individual are safeguarded. This foundational approach challenges traditional views of justice that often prioritize the interests of the powerful or privileged.
\\ \textit{Note:} correct because it accurately describes how Rawls's veil of ignorance creates a fair decision-making environment by removing personal biases and self-interest, prompting individuals to prioritize justice and fairness in their choices.
\item \textbf{Answer:} The Accident fallacy occurs in philosophy when a generalization is improperly applied to a specific case, leading to erroneous conclusions. This fallacy arises when an essential characteristic of the generalization is overlooked, resulting in misapplication. For example, consider the generalization that "all birds can fly." If one were to argue that since penguins are birds, they must be able to fly, this would be a fallacy, as it ignores the specific characteristics of penguins. Such misapplications can distort reasoning and hinder accurate understanding of the nuances within categories, highlighting the importance of context when applying general principles.
\\ \textit{Note:} correct because it clearly defines the Accident fallacy as the improper application of a generalization to a specific case, illustrated with the example of penguins, which highlights the oversight of essential characteristics that leads to erroneous conclusions.
\end{enumerate}

\vfill
\noindent\textit{End of Answer Key}
\end{document}